% sections/introduction.tex
% Phase 9: Introduction (RAPP-01)

\section{Introduction}
\label{sec:introduction}

Le 5 ao\^ut 2024, les march\'es financiers mondiaux ont c\'ed\'e en quelques
heures, un \'episode que la presse a appel\'e le \og Lundi Noir \fg{}. Le
d\'eclencheur est mon\'etaire~: le 31 juillet, la Banque du Japon rel\`eve ses
taux directeurs par surprise~\cite{boj2024rate}, ce qui force le d\'enouement
des positions de \textit{carry trade} sur le yen. Le m\'ecanisme est connu~:
des investisseurs empruntent en yen \`a taux quasi nul (le Japon maintient des
taux n\'egatifs depuis 2016) pour placer ailleurs \`a rendement plus
\'elev\'e~\cite{reuters2024carrytrade}. Quand le yen s'appr\'ecie brutalement,
tout le monde doit d\'eboucler en m\^eme temps pour couvrir les pertes de
change. Les ventes forc\'ees se r\'epercutent sur toutes les classes d'actifs,
y compris la crypto~: sur Polygon, blockchain de couche~2 d'Ethereum, les
protocoles de finance d\'ecentralis\'ee (DeFi) encaissent des vagues de
liquidations et de transferts de panique. Sur la seule journ\'ee du 5 ao\^ut,
plus de 401\,000 transferts du jeton WETH (Wrapped Ether) sont
enregistr\'es sur Polygon.

Ces transferts forment un r\'eseau orient\'e pond\'er\'e. Chaque adresse
Ethereum est un sommet, chaque transfert une ar\^ete orient\'ee de l'\'emetteur
vers le r\'ecepteur, et le montant transf\'er\'e d\'efinit le poids de l'ar\^ete.
La th\'eorie des graphes donne des outils pour analyser ce r\'eseau~:
topologie, acteurs centraux, structure communautaire, dynamique temporelle
des \'echanges~\cite{newman2018networks}. C'est plus utile qu'un simple
d\'ecompte de transactions, parce qu'on peut alors voir o\`u se concentrent
les flux, quels protocoles servent de relais, et comment le r\'eseau se
fragmente quand le march\'e tangue.

Le rapport s'organise autour de six questions de recherche (RQ)~:

\begin{enumerate}
    \item \textbf{RQ1 --- Topologie du r\'eseau~:} analyser la structure globale
    \`a travers la distribution des degr\'es, la densit\'e, les composantes connexes
    et le diam\`etre du graphe.
    \item \textbf{RQ2 --- Centralit\'e~:} identifier les n\oe{}uds les plus
    influents par quatre mesures de centralit\'e (degr\'e, \textit{betweenness},
    \textit{closeness}, PageRank) et interpr\'eter leurs r\^oles.
    \item \textbf{RQ3 --- Communaut\'es~:} d\'etecter la structure communautaire
    par l'algorithme de Louvain, mesurer la modularit\'e et analyser la
    distribution des tailles de communaut\'es.
    \item \textbf{RQ4 --- Dynamique temporelle~:} d\'ecouper la journ\'ee en
    fen\^etres horaires et suivre l'\'evolution des m\'etriques du r\'eseau pour
    identifier les pics d'activit\'e.
    \item \textbf{RQ5 --- Propri\'et\'es petit-monde~:} \'evaluer le coefficient
    de clustering et la longueur moyenne des chemins, puis comparer au mod\`ele
    al\'eatoire d'Erd\H{o}s-R\'enyi.
    \item \textbf{RQ6 --- Flux pond\'er\'es~:} analyser la distribution des
    volumes WETH \'echang\'es, quantifier la concentration par le coefficient de
    Gini et la courbe de Lorenz.
\end{enumerate}

L'environnement de travail et le jeu de donn\'ees sont d\'ecrits en premier,
puis les six questions sont trait\'ees une par une (m\'ethodologie,
r\'esultats, interpr\'etation). Le bilan revient sur ce que les six analyses
disent ensemble, et propose quelques pistes pour aller plus loin.

En compl\'ement, un prototype web ouvre les m\^emes donn\'ees \`a
l'exploration directe \`a la souris~: \url{https://crypto-graph-explorer.vercel.app}~\cite{cryptographexplorer_web}.
On peut y zoomer dans le graphe, filtrer par communaut\'e, ouvrir un
n\oe{}ud, et relancer les six RQ et les cinq analyses compl\'ementaires
(section~\ref{sec:complementaires}) sur des m\'etriques pr\'e-calcul\'ees.
Les chiffres exacts rapport\'es ici viennent de l'ex\'ecution Python
d\'ecrite dans l'annexe~\ref{ann:manuel}.
